\section{Diseño metodológico}

\subsection{Estrategia general de solución}

La estrategia general de solución consiste en modelar la malla curricular como un Grafo Dirigido Acíclico (DAG), en el que los nodos representan asignaturas y las aristas codifican las relaciones de prerrequisito, para luego formular la planificación académica como un problema de búsqueda y optimización bajo restricciones sobre dicho grafo. A partir de esta representación, el sistema genera configuraciones semestrales que respetan el orden topológico del DAG, las restricciones de carga académica y las políticas institucionales, evaluando cada configuración mediante una función de costo que sintetiza el desempeño de la ruta propuesta.

Metodológicamente, el proceso se desarrolla en cuatro pasos principales. En primer lugar, se construye la base de conocimiento curricular a partir de la información institucional: asignaturas, créditos, prerrequisitos y oferta por periodo. En segundo lugar, se define el modelo del estudiante, incorporando su historial académico y sus preferencias (como carga máxima por semestre o énfasis deseado). En tercer lugar, el motor de planificación explora el espacio de rutas posibles mediante algoritmos de grafos y heurísticas de generación de planes, descartando aquellas que violan restricciones estructurales o de carga. Finalmente, se calculan métricas de desempeño para cada ruta viable y se selecciona la que mejor satisface los criterios definidos en el marco PEAS.

Esta estrategia se justifica porque permite integrar en un mismo esquema la representación explícita del conocimiento curricular y un mecanismo de razonamiento que combina reglas duras (prerrequisitos, límites de créditos) con criterios de optimización (tiempo de graduación, distribución de cursos críticos). Además, el uso de un DAG y de algoritmos estándar de grafos facilita la extensibilidad del sistema, permitiendo incorporar en fases posteriores modelos de predicción de rendimiento o riesgo de reprobación sin modificar la estructura metodológica fundamental.

\pagebreak

